Дано целое число $A$. Проверить истинность высказывания: «Число $A$ является положительным и двузначным».
Дано целое число $A$. Проверить истинность высказывания: «Число $A$ является нечетным и двузначным».
Дано целое число $A$. Проверить истинность высказывания: «Число $A$ является четным и двузначным».
Даны два целых числа: $A$, $B$. Проверить истинность высказывания: «Справедливы неравенства $A > 2$ и $B \leq 3$».
Даны два целых числа: $A$, $B$. Проверить истинность высказывания: «Справедливы неравенства $A \geq 0$ или $B < -2$».
Даны три целых числа: $A$, $B$, $C$. Проверить истинность высказывания: «Справедливо двойное неравенство $A < B < C$».
Даны три целых числа: $A$, $B$, $C$. Проверить истинность высказывания: «Число $B$ находится между числами $A$ и $C$».
Даны два целых числа: $A$, $B$. Проверить истинность высказывания: «Каждое из чисел $A$ и $B$ нечетное».
Даны два целых числа: $A$, $B$. Проверить истинность высказывания: «Хотя бы одно из чисел $A$ и $B$ нечетное».
Даны два целых числа: $A$, $B$. Проверить истинность высказывания: «Ровно одно из чисел $A$ и $B$ нечетное».
Даны два целых числа: $A$, $B$. Проверить истинность высказывания: «Числа $A$ и $B$ имеют одинаковую четность».
Даны три целых числа: $A$, $B$, $C$. Проверить истинность высказывания: «Каждое из чисел $A$, $B$, $C$ положительное».
Даны три целых числа: $A$, $B$, $C$. Проверить истинность высказывания: «Хотя бы одно из чисел $A$, $B$, $C$ положительное».
Даны три целых числа: $A$, $B$, $C$. Проверить истинность высказывания: «Ровно одно из чисел $A$, $B$, $C$ положительное».
Даны три целых числа: $A$, $B$, $C$. Проверить истинность высказывания: «Ровно два из чисел $A$, $B$, $C$ являются положительными».
Дано целое положительное число. Проверить истинность высказывания: «Данное число является четным трехзначным».
Дано целое положительное число. Проверить истинность высказывания: «Данное число является нечетным трехзначным».
Проверить истинность высказывания: «Среди трех данных целых чисел есть хотя бы одна пара совпадающих».
Проверить истинность высказывания: «Среди трех данных целых чисел есть хотя бы одна пара взаимно противоположных».
Дано трехзначное число. Проверить истинность высказывания: «Все цифры данного числа различны».
Даны числа $x$, $y$. Проверить истинность высказывания: «Точка с координатами $(x, y)$ лежит во второй координатной четверти».
Даны числа $x$, $y$. Проверить истинность высказывания: «Точка с координатами $(x, y)$ лежит в четвертой координатной четверти».
Даны числа $x$, $y$. Проверить истинность высказывания: «Точка с координатами $(x, y)$ лежит во второй или третьей координатной четверти».
Даны числа $x$, $y$. Проверить истинность высказывания: «Точка с координатами $(x, y)$ лежит в первой или третьей координатной четверти».
Даны числа $x$, $y$, $x_1$, $y_1$, $x_2$, $y_2$. Проверить истинность высказывания: «Точка с координатами $(x, y)$ лежит внутри прямоугольника, левая верхняя вершина которого имеет координаты $(x_1, y_1)$, правая нижняя --- $(x_2, y_2)$, а стороны параллельны координатным осям».
Дано трехзначное число. Проверить истинность высказывания: «В числе есть одинаковые цифры».
Дано четырехзначное число. Проверить истинность высказывания: «В числе первая и вторая цифры одинаковы».
Дано четырехзначное число. Проверить истинность высказывания: «В числе вторая и третья цифры одинаковы».
Дано четырехзначное число. Проверить истинность высказывания: «В числе первая и вторая цифры не совпадают с последней».
Дано четырехзначное число. Проверить истинность высказывания: «В числе первая цифра совпадает с последней».
